\chapter{Project Overview}\label{chap:intro}

For the final week of the Ejada Cloud Build Track internship, our Group A team
split an e-commerce application into five independently deployable Git
repositories: four backend services and a Next.js frontend. We then deployed the
system to Oracle Kubernetes Engine (OKE) using a Terraform-provisioned cluster
and a Helm/ArgoCD GitOps pipeline.

\section{Objectives and Scope}\label{sec:core-question}
We set two main objectives for the build: first, to split the application into
microservice boundaries based on data ownership, each with its own repository and
container image; and second, to automate the path from a code push to a running
pod.

\textbf{For our initial scope, we focused on:} one OCI region (\texttt{me-jeddah-1}), one compartment
(\texttt{shared-group-a-cmp}), one OKE cluster and namespace
(\texttt{ecommerce-ns}), the five application repositories above, and the
\texttt{Terraform-code} repository holding the infrastructure, Kubernetes manifests,
and Helm chart.

\textbf{We initially treated as out of scope:} automated testing, multi-region or
multi-availability-domain high availability, and a production security review.
However, we later added a dedicated QA and remediation cycle
(Section~\ref{sec:testing}). We also found that three customer-facing storefront
pages remain incomplete (Section~\ref{sec:frontend-completeness}) and that a few
areas of our implementation differ from the original design draft, as we
summarise in Section~\ref{sec:draft-vs-implementation}.

\section{The Journey from Monolith to Microservices}\label{sec:monolith-journey}
The project began as a single monolithic repository (\textit{as-is}) containing both the Next.js frontend and a unified Express.js backend. This architecture presented significant scaling and deployment bottlenecks. Because the entire system was tightly coupled, a single code change—whether a minor UI tweak or a critical payment logic fix—necessitated rebuilding and redeploying the entire application stack. This lack of isolation also meant that a failure in one domain (e.g., product catalog search) could potentially crash the entire backend, degrading the user experience across all functionalities.

To achieve greater fault tolerance, team autonomy, and independent scalability, our team embarked on a comprehensive refactoring journey to transition the application into a microservices architecture. We decoupled the monolith based on bounded business domains into four distinct backend services: \texttt{auth-service}, \texttt{products-service}, \texttt{orders-service}, and \texttt{payments-service}, alongside the standalone \texttt{ecomm-ui} Next.js frontend (acting as a Backend-For-Frontend). 

Each microservice was granted its own dedicated Git repository, independent Dockerfile, and distinct CI/CD workflow, culminating in our current multiple-repo configuration. This separation of concerns allows individual teams to develop, test, and ship features independently. For instance, the orders service can scale independently of the auth service during high-traffic checkout events. Furthermore, we containerized these services and deployed them to Oracle Kubernetes Engine (OKE) using a Terraform $\rightarrow$ Docker $\rightarrow$ Helm $\rightarrow$ ArgoCD pipeline, enabling automated GitOps-driven deployments. While the services are now independently deployable, they currently share a single MongoDB instance, representing a pragmatic intermediate step toward full data-tier isolation.

\section{Project Ownership}\label{sec:key-stakeholders}
We summarise each team member's area of responsibility in
Table~\ref{tab:stakeholder-map}.

\begin{table}[H]
\centering
\caption{Team members and areas of responsibility.}
\label{tab:stakeholder-map}
\small
\renewcommand{\arraystretch}{1.25}
\begin{tabularx}{\textwidth}{@{}
                >{\raggedright\arraybackslash}p{4.3cm} >{\raggedright\arraybackslash}X
                @{\hspace{1.8em}}
                >{\raggedright\arraybackslash}p{3.9cm} >{\raggedright\arraybackslash}X @{}}
\toprule
\textbf{Member} & \textbf{Area} & \textbf{Member} & \textbf{Area} \\
\midrule
Ali Ahmed Fakhry Hamad & Deployment       & Mohamed Kholosy  & E-commerce application \\
Abdelrahman Darwish    & E-commerce application & Habiba Fawzy & Database \\
Salma Azara            & Terraform code   & Chris Mina Kamal & Terraform code \\
Ali Tallawy            & Testing          & Basel Alaa       & Terraform code \\
Beshoy Sorial          & CI/CD            & Yara Adel        & Architecture and design \\
Ali Youssef            & Documentation    & Peter Kameel     & Testing \\
\bottomrule
\end{tabularx}
\renewcommand{\arraystretch}{1}
\end{table}
